\subsection{Beantwortung der Fragestellungen}
\label{subsec:beantwortung-fragestellungen}

\paragraph{Fragestellung 1: Adressierte Anforderungen.}
Das Technologie-Template adressiert die in Kapitel~\ref{sec:anforderungen}
formulierten fachlich-technischen, qualitativen und betrieblichen Anforderungen
durch eine generierbare Baseline, klare Komponentenstrukturen,
profilabhängige Auswahl, gemeinsame Workflows, zentrale Pflege sowie
versionierte Konfiguration und Capability-Abhängigkeiten. Getrennte
Laufzeitgrenzen und profilabhängige Prüfpfade stützen dies. Fachlogik,
Authentifizierung, Berechtigungen und Produktivbetrieb bleiben
produktspezifisch.

\paragraph{Fragestellung 2: Varianten für Web, Cloud und Desktop.}
Architektur und Profilmodell kombinieren den gemeinsamen Frontend-Kern mit
optionalen Backend-, Cloud- und Tauri-Bausteinen. Fünf Presets bilden die
Plattformzuschnitte ab; PostgreSQL wird getrennt und nur bei Backendfähigkeit
aufgelöst. Generierte Projekte enthalten so überwiegend nur benötigte Pfade.
Dies gilt für den untersuchten Katalog; \texttt{desktop-cloud} und
\texttt{full-platform} sind technisch gleich gescaffoldet und vor allem
semantisch abgegrenzt.

\paragraph{Fragestellung 3: Wiederverwendung.}
Gemeinsamer Kern und Single-Repository-Strategie fördern Wiederverwendung und
Standardisierung weitgehend, weil Implementierungen, Profile, Generator,
Konfiguration, Tests und Dokumentation zentral gepflegt werden. Deklarative
Presets und gemeinsame Tooling-Befehle begrenzen abweichende Ausgangsstände.
Dieser Vorteil gilt unmittelbar für die Baseline und künftige Generierungen;
bereits abgeleitete Produktprojekte bleiben separate Repositories ohne
automatischen Rückfluss späterer Template-Änderungen.

\paragraph{Fragestellung 4: Reduktion wiederkehrender Aufwände.}
Für Projektinitialisierung, Entwicklungsworkflow, Qualitätssicherung und
Bereitstellung besteht begründbares Einsparpotenzial: Profilwahl, Generierung,
Diagnose, Installation, Start, Tests und Builds sind automatisiert oder
vereinheitlicht und teilweise technisch abgenommen. Messdaten zu Zeit und
Aufwand sowie Vergleichsdaten fehlen. Eine Produktivitätssteigerung ist daher
nicht nachgewiesen; zudem sind zentrale Pflege und mehrere Toolchains
gegenzurechnen.

\paragraph{Fragestellung 5: Projekteignung.}
Eine geeignete Ausgangsbasis stellt das Template grundsätzlich für technisch
verwandte Web-, Desktop- und kombinierte Business-Anwendungen dar, deren
Plattform-, Cloud- und Persistenzbedarf abbildbar ist. Stark native Anwendungen,
Spiele und spezialisierte Echtzeitsysteme liegen nicht im primären Zielbereich.
Kapitel~\ref{sec:einsatz-transfer} leitet eine technikneutrale Analyse, das
kleinste passende Profil und die getrennte Capability-Auswahl ab. Die
Einstufung bleibt qualitativ und ist keine universelle Eignungszusage.

\paragraph{Fragestellung 6: Grenzen, Wartungsaufwände und Risiken.}
Zu berücksichtigen sind die wachsende Regressionsbreite des Master-Templates,
manuell synchronisierte Schnittstellen, die nicht strikt getrennte
Capability-Auswahl und fehlende automatische Updates generierter Projekte.
Mehrere Toolchains erhöhen den Pflegebedarf. Aktuell schlagen Tooling-,
Desktopprofil- und Windows-CI-Pfade fehl; Release Validation, Compose-Start
und Branch Protection sind nicht verifiziert. Signing, Notarisierung, Updater,
Produktiv-Deployment und Sicherheitsmodell bleiben Produktaufgaben.
